\section{Transfer Validation of the Chart-CNN Paradigm}

Before comparing image and tabular representations, we first validate that the chart-CNN paradigm of \citet{jiang2023reimagining} transfers to China A-shares. This validation is not a claim that the CNN dominates all alternatives. It is a narrower requirement: the migrated image model should preserve the qualitative predictability patterns of the original design.

Our implementation aligns with the original paradigm along four dimensions: OHLCV histories are rendered as grayscale price charts; the CNN architecture follows the image-window-dependent convolutional design; the main experiments use five-model ensembles; and the I20/R5 specification, a 20-day image predicting roughly five-day returns, is treated as the principal chart-image specification.

\begin{table}[htbp]
\centering
\caption{Transfer validation of Jiang et al. chart-CNN specifications}
\label{tab:transfer_validation}
\begin{tabular}{lcrr}
\toprule
Specification & Image window & EW Sharpe & VW Sharpe \\
\midrule
I20/R5 & 20 days & 2.39 & 1.01 \\
\bottomrule
\end{tabular}
\begin{minipage}{0.92\linewidth}
\footnotesize
\textit{Notes:} Sharpe uses arithmetic mean divided by period volatility, annualized by $\sqrt{52}$. This transfer-validation row uses the original transfer sample; subsequent matched-model comparisons use the common OOS panel. Compounded annual return is reported separately where available.
\end{minipage}
\end{table}


Table~\ref{tab:transfer_validation} shows that the principal transfer specification is empirically meaningful. Under the common arithmetic Sharpe convention, I20/R5 reaches an equal-weighted high-minus-low Sharpe of 2.39 and a value-weighted Sharpe of 1.01. Historical alternative transfer variants are not included because their portfolio summaries have not been regenerated under this common convention.

This evidence validates the migration in the sense needed for this paper. The chart-image CNN extracts economically meaningful return-predictive information in China A-shares. The remaining question is therefore whether its predictive content is unique relative to nonlinear tabular learners using the same OHLCV information.
